

\textbf{\SGAInlineRefTarget{sga3:xxiv:lemma:4-4-4}Lemma 4.4.4.} \emph{Let \(S\) be a scheme, \(G\) and \(H\) two \(S\)-group schemes, \(G \to \operatorname{Aut}_{S-gr.}(H)\) an
action of \(G\) on \(H\), \(E\) a \(G\)-principal homogeneous bundle, \(F\) an \(H_0\)-principal homogeneous bundle, where \(H_0\) is
the form of \(H\) associated to \(E\). Then \(E \times_{S} F\) is endowed naturally with a structure of principal homogeneous
bundle under the semi-direct product \(H \cdot G\).}

Write \((e, g) \mapsto eg\) (resp. \((f, u) \mapsto f u\)) for the action (on the right) of \(G\) on \(E\) (resp. of \(H_0\) on
\(F\)). Write

\[
  p:E\times_S H\longrightarrow H_0.
\]

for the canonical projection (\(H_0\) is by definition the quotient of \(E \times_{S} H\) by \(G\) acting on it by the formula
\((e, h) g = (eg, g^{-1} h g)\)). Consider the morphism

\[
  r:E\times_S F\times_S H\times_S G\longrightarrow E\times_S F.
\]

defined set-theoretically by

\[
  r(e,f,h,g)=\bigl(eg,f\cdot p(e,h)\bigr).
\]

The morphism \(r\) indeed defines an action of the semi-direct product \(H \cdot G\) on \(E \times_{S} F\). Indeed, one has
set-theoretically

\[
  r\bigl(r(e,f,h,g),h',g'\bigr)
  =\bigl(egg',f\,p(e,h)\,p(eg,h')\bigr).
\]

but

\[
  p(e,h)p(eg,h')
  =p(e,h)p(e,gh'g^{-1})
  =p(e,hgh'g^{-1}).
\]

whence

\[
  r\bigl(r(e,f,h,g),h',g'\bigr)
  =r(e,f,hgh'g^{-1},gg').
\]

which had to be proved.

To prove now that this law is indeed a law of principal homogeneous bundle, one may assume that \(E\) and \(F\) are trivial,
in which case one sees at once that \(E \times_{S} F\) is also a trivial bundle.

\subsection*{\SGAInlineRefTarget{sga3:xxiv:subsection:4-5}4.5. Isotriviality of maximal tori of semisimple groups}
\footnote{\textit{Editors' note.} This no.~\SGARefLink{sga3:xxiv:subsection:4-5}{4.5} has been added.}

Let us bring out here the following result, contained implicitly in \SGARefLink{sga3:xxiv:2.6}{2.6} (cf.~the N.D.E. (13)) and which was mentioned at
the end of \SGARefLink{sga3:ix:comments:1:2}{Exp. IX 1.2}. Recall (\SGARefLink{sga3:XXIII:5.11}{XXIII 5.11}) that, for every scheme \(S\) and every reduced root datum \(R\), one writes
\(\acute{E}p_{S}(R)\) for the unique pinned \(S\)-group of type \(R\) (which exists by Exp. XXV) and \(T_{S}(R)\) its canonical
maximal torus.

\textbf{\SGAInlineRefTarget{sga3:xxiv:proposition:4-5-1}Proposition 4.5.1.} \emph{Let \(G\) be an \(S\)-semisimple group of constant type \(R\); assume that \(G\) possesses a maximal
torus \(T\) (which is the case if \(S\) is semi-local, by \SGARefLink{sga3:xiv:3.20}{XIV 3.20}).}

\emph{\SGAInlineRefTarget{sga3:xxiv:proposition:4-5-1:item:i}(i) Then \(P = Isom_{S-gr.}(\acute{E}p_{S}(R), T_{S}(R); G, T) / T_{S}(R)_{ad}\) is a principal \(S\)-bundle under the
finite group \(\operatorname{Aut}(R)\), and \(T_P\) is a split maximal torus of \(G_P\). Consequently, \(T\) is isotrivial.}

\emph{\SGAInlineRefTarget{sga3:xxiv:proposition:4-5-1:item:ii}(ii) Moreover, for every finite set \(F\) of points of \(P\) contained in an affine open set, there exists an open set \(U\)
of \(P\) containing \(F\) such that \(G_U\) is split.}

\emph{Proof.} The exact sequence

\[
  1\longrightarrow T_S(\mathscr R)_{\mathrm{ad}}
  \longrightarrow
  \underline{\operatorname{Aut}}_{S\text{-}\mathrm{gr.}}
  \bigl(\operatorname{\acute{E}p}_S(\mathscr R),T_S(\mathscr R)\bigr)
  \longrightarrow
  \underline{\operatorname{Aut}}_{S\text{-}\mathrm{gr.}}\bigl(T_S(\mathscr R)\bigr).
\]

induces for every \(S' \to S\) a morphism of bundles \(P_{S'} \to Isom_{S'-gr.}(T_{S'}(R), T_{S'})\), so if \(P_{S'}\)
possesses a section, then \(T_{S'}(R)\) and \(T_{S'}\) are isomorphic. This is in particular the case for \(S' = P\).

Note moreover that \(\operatorname{Aut}_{S-gr.}(\acute{E}p_{S}(R), T_{S}(R)) / T_{S}(R)_{ad}\) is none other than the
constant \(S\)-group \(\operatorname{Aut}(R)_{S}\), and that \(P\) is identified with the \(\operatorname{Aut}(R)_{S}\)-torsor
\(Isom(R, \tilde{R}(G, T))\), where \(\tilde{R}(G, T)\) denotes the twisted root datum associated to \((G, T)\) (cf.~\SGARefLink{sga3:xxii:prop:1.10}{XXII
1.10}).

Then \SGARefLink{sga3:xxiv:proposition:4-5-1:item:ii}{(ii)} follows from \SGARefLink{sga3:xxiv:proposition:4-5-1:item:i}{(i)} and from \SGARefLink{sga3:xxiv:2.14}{2.14}, taking into account the fact that the \(g_{\alpha}\) (\(\alpha\) ranging over the
roots of \(R\)) are free \(\mathcal O_U\)-modules of rank 1 on a sufficiently small affine open set \(U\) containing \(F\).

\section*{\SGAInlineRefTarget{sga3:xxiv:section:5}5. Canonical decomposition of an adjoint or simply connected group}

In this section, we shall use the results of no 1 to generalize to the case of schemes a classical decomposition of
adjoint (resp. simply connected) groups. So as not to overburden the exposition indefinitely, the proofs are sketched
and the details are left to the reader; in fact it is always a matter of absolutely standard proofs in the theory of
principal bundles: reduction of the structure group, twisting, etc.

\subsection*{\SGAInlineRefTarget{sga3:xxiv:subsection:5-1}5.1.}

Recall (\SGARefLink{sga3:XXI:7.4}{Exp. XXI, 7.4}) that a Dynkin diagram is a disjoint union of its connected components, which are Dynkin diagrams.
Moreover, every non-empty connected Dynkin diagram corresponding to a root datum is isomorphic to one of the standard
diagrams (\(A_{n}\), \(B_{n}\), \ldots, \(G_2\)) which were exhibited in \SGARefLink{sga3:XXI:7.4.6}{Exp. XXI, 7.4.6}. In the sequel, we shall be interested
only in Dynkin diagrams whose connected components are of one of the preceding types. Let \(T\) be the set of these
standard diagrams. For every Dynkin diagram \(D\), let \(n(t) = n_{D}(t)\) be the number of connected components of \(D\)
isomorphic to \(t\), where \(t \in T\). The \emph{type} of \(D\) is by definition \(\sum_{t \in T} n_{D}(t) \cdot t\).

A Dynkin diagram of type \(t\) is said to be \emph{simple of type \(t\)}, a Dynkin diagram of type \(n \cdot t\) is said to be
\emph{isotypic of type \(t\)}. Let \(D_0\) be the set of connected components of \(D\) and let

\[ a : D_{0} \to T \]

be the obvious map. The type of \(D\) is none other than \(\sum_{x \in D_{0}} a(x)\).

\subsection*{\SGAInlineRefTarget{sga3:xxiv:subsection:5-2}5.2.}

Let \(S\) be a scheme, \(D\) an \(S\)-Dynkin scheme (verifying the restrictive condition stated above). The cokernel of the
pair of morphisms \(L \Rightarrow D\) (\(L =\)\footnote{N.D.E.: We have replaced \(R\) by \(L\), as in \SGARefLink{sga3:xxiv:3.2}{3.2}.} scheme of bonds of \(D\)) is written \(D_0\). This is the ``scheme
of connected components'' of \(D\) (it

exists trivially when \(D\) is constant; the general case is deduced from this by descent); it is a finite twisted
constant \(S\)-scheme. One has a canonical morphism

\[ a : D_{0} \to T_{S}. \]

For \(t \in T\), set \(a^{-1}(t) = D_{0, t}\); this is a subscheme of \(D_0\), whose inverse image in \(D\), written \(D_{t}\), is
the \emph{isotypic component of type \(t\)} of the Dynkin scheme \(D\). Each \(D_{t}\) is a subscheme of \(D\), and one has

\[
  D=\coprod_{t\in T}D_t.
\]

Note that the degree of \(D_{0, t}\) at \(s \in S\) is \(n(s, t)\), if the type of \(D\) at \(s\) is \(\sum_{t} n(s, t) \cdot t\).

\subsection*{\SGAInlineRefTarget{sga3:xxiv:subsection:5-3}5.3.}

In what follows, we shall consider only semisimple adjoint (resp. simply connected) groups. To simplify the language, we
shall state the results for adjoint groups; all statements will remain valid if one substitutes everywhere ``simply
connected'' for ``adjoint''.

Recall that an adjoint reduced root datum is determined up to isomorphism by the type of its Dynkin diagrams. We shall
therefore say that an adjoint root datum \(R\) (resp. a semisimple adjoint group \(G\)) is \emph{of type \(\sum n(t) \cdot t\)} if
its Dynkin diagrams are (resp. if its type is given by an adjoint root datum of type \(\sum n(t) \cdot t\)). We shall say
that \(R\) or \(G\) is \emph{simple of type \(t\)} (resp. \emph{isotypic of type \(t\)}) if its type is \(t\) (resp. \(n \cdot t\), \(n > 0\)).

If \(G\) is a semisimple adjoint group, we shall use the symbols \(Dyn_{0}(G)\) and \(Dyn_{0, t}(G)\) in the sense defined in
\SGARefLink{sga3:xxiv:subsection:5-2}{5.2}.

\subsection*{\SGAInlineRefTarget{sga3:xxiv:subsection:5-4}5.4.}

Let \(t_i\), \(i=1,2,\ldots,n\), be distinct elements of \(T\), and let \(R_i\) be a pinned adjoint root datum isotypic of
type \(t_{i}\). Consider the product \(R = R_{1} \times \cdot\cdot\cdot \times R_{n}\) (\SGARefLink{sga3:XXI:6.4.1}{Exp. XXI, 6.4.1}). Let \(S\) be a
scheme such that the different \(\acute{E}p_{S}(R_{i})\) exist (cf.~Exp. XXV). One verifies at once that there exists a
canonical isomorphism

\[
  \tag{*}\label{sga3:xxiv:subsection:5-4:equation:star}
  \operatorname{\acute{E}p}_S(\mathscr R)
  =
  \operatorname{\acute{E}p}_S(\mathscr R_1)
  \times_S\cdots\times_S
  \operatorname{\acute{E}p}_S(\mathscr R_n).
\]


